\documentclass{article}
\usepackage[a4paper, total={6.5in, 10in}]{geometry}
\setlength{\parindent}{0pt}
\usepackage{tcolorbox}
\tcbset{colback=white!94!black, colframe=black!60!white, coltitle = white, boxrule=0.8pt, arc=2mm}
\newtcolorbox{boxs}[2][]{title= #2,#1}
\usepackage{datetime2}
\usepackage[british]{babel}
\usepackage{amsfonts}
\usepackage{amsmath}

\title{\textbf{Linear Algebra HW 6}}
\author{Robert O'Connell}
\date{\today}
\begin{document}
\maketitle
\vspace{5mm}
\subsection*{Question 1}
If the image of a linear transformation is equal to the kernel of the same linear transformation, then by the
rank-nullity theorem we can see that \(\dim(\mathrm{Im} \ T) = \dim(\ker T) = 2\)
\\

i.e. we need to send a two dimensional subspace to \(\vec{0}\) (\(\ker T\)) and then the remaining two dimensional needs to be mapped to 
the the \(\ker T\)

One such tranformation is \(T: \mathbb{R}^4 \rightarrow \mathbb{R}^4\), \(T(\vec{v}) = A \vec{v} \quad \forall \ \vec{v} \in \mathbb{R}^4\)
\[
A = \begin{bmatrix}
    0 & 0 & 1 & 0 \\
    0 & 0 & 0 & 1 \\
    0 & 0 & 0 & 0 \\
    0 & 0 & 0 & 0 
\end{bmatrix}
\]
As the first  two basis vectors get mapped to \(\vec{0}\) and the other two basis vectors get mapped to the first two basis vectors

\subsection*{Question 2}
Let \(\vec{v}_1, \vec{v}_2 \in V\) and \(\alpha \in \mathbb{R}\),
\begin{itemize}
    \item \((S+T)(\vec{v}_1+ \vec{v}_2) = S(\vec{v}_1+ \vec{v}_2) + T(\vec{v}_1+ \vec{v}_2) = S\vec{v}_1+ T\vec{v}_1+ S\vec{v}_2 + T\vec{v}_2 = (S+T)(\vec{v}_1) + (S+T)(\vec{v}_2)\)
    \item \((S+T)(\alpha \vec{v}_1) = S(\alpha \vec{v}_1) + T(\alpha \vec{v}_1) = \alpha(S(\vec{v}_1) + T(\vec{v}_1))\)
    \item Thus, \((S+T)\) is a linear transformation
    \\
    \item \((\lambda T)(\vec{v}_1 + \vec{v_2}) = \lambda((T)(\vec{v}_1+ \vec{v}_2)) =\lambda(T(\vec{v}_1)+ T(\vec{v}_2)) = \lambda T(\vec{v}_1)+ \lambda T(\vec{v}_2)\)
    \item \((\lambda T)(\alpha\vec{v}_1) = \lambda(T(\alpha \vec{v}_1)) = \lambda\alpha T(\vec{v}_1) = \alpha(\lambda T(\vec{v}_1))\)
    \item Thus \((\lambda T)\) is a linear transformation
\end{itemize}

\subsection*{Question 3}
\(A\) diagonalizable \(\iff\) \(\exists\) a basis of \(\mathbb{R}^3\) consisting of eigenvectors of \(A\) \(\Leftarrow\) \(A\) has \(n\) distinct eigenvalues

To first find the eigenvalues we solve the characteristic polynomial of \(A\) equal to zero

\[
\chi_A(\lambda) = \det(A-\lambda I) = 0
\]
\begin{align*}
    \det(A-\lambda I) &= \begin{bmatrix}
        6 - \lambda & -10 & 1 \\
        1 & -11 -\lambda & 3 \\
        4 & -40 & 11 - \lambda
    \end{bmatrix}
    \\
    &= (6-\lambda)((-11-\lambda)(11 - \lambda) + 120) + (10)((11 - \lambda) -12) + (-40 + 4(11 + \lambda))
    \\
    &= (6 - \lambda)(\lambda^2 -1) -10(\lambda + 1) + 4 + 4\lambda
    \\
    &= (6- \lambda)(\lambda + 1)(\lambda -1) - 10(\lambda + 1) + 4(\lambda +1)
    \\
    &= (\lambda+1)(6\lambda - 6 - \lambda^2 + \lambda -10 +4)
    \\
    &= -(\lambda + 1)(\lambda -3)(\lambda -4)
\end{align*}

We have three distinct eigenvalues \(\lambda = -1, \lambda = 3\) and \(\lambda = 4\),

Hence \(A\) is diagonalizable
\\

To find the diagonal matrix \(D\) to which is \(A\) is similar we note that in the basis of \(\mathbb{R}^3\) consisting 
of eigenvectors, say \(\mathcal{B} = \{\vec{b}_1, \vec{b}_2, \vec{b_3}\}\), that each of them is only scaled, hence \([A]_{\mathcal{B},\mathcal{B}}\) is diagonal.
\[
[A]_{\mathcal{B}, \mathcal{B}} = \begin{bmatrix}
    -1 & 0 & 0 \\
    0 & 3 & 0 \\
    0 & 0 & 4
\end{bmatrix}
\]

Thus \(D = B^{-1} A B\) where \(D = [A]_{\mathcal{B}, \mathcal{B}}\) and \(B = \mathcal{P}_{\mathcal{B}\to \mathcal{E}}\)
\\

To find the eigenvectors we look in \(\mathcal{N}(A-\lambda I)\),
\begin{align*}
    \mathcal{N}(A + I) &, \begin{bmatrix}
        7 & -10 & 1 &| & 0 \\
        1 & -10 & 3 &|& 0 \\
        4 & -40 & 12& |&  0
    \end{bmatrix}
    \\
    &= \begin{bmatrix}
        3 & 0 & -1 &|&0 \\
        0 & 3 & -1 &|&0 \\
        0 & 0 & 0 &|&0
    \end{bmatrix}
\end{align*}
One such eigenvector is \(\begin{bmatrix}
    1 \\ 1 \\ 3
\end{bmatrix}\)

\begin{align*}
    \mathcal{N}(A-3I) &, \begin{bmatrix}
        3 & -10 & 1 &| & 0 \\
        1 & -14 & 3 &|& 0 \\
        4 & -40 & 8& |&  0
    \end{bmatrix}
    \\
    &= \begin{bmatrix}
        4 & 0 & -2 &|&0 \\
        0 & 4 & -1 &|&0 \\
        0 & 0 & 0 &|&0
    \end{bmatrix}
\end{align*}
One such eigenvector is \(\begin{bmatrix}
    2 \\ 1 \\ 4
\end{bmatrix}\)

\begin{align*}
    \mathcal{N}(A-3I) &, \begin{bmatrix}
        2 & -10 & 1 &| & 0 \\
        1 & -15 & 3 &|& 0 \\
        4 & -40 & 7& |&  0
    \end{bmatrix}
    \\
    &= \begin{bmatrix}
        4 & 0 & -3 &|&0 \\
        0 & 4 & -1 &|&0 \\
        0 & 0 & 0 &|&0
    \end{bmatrix}
\end{align*}
One such eigenvector is \(\begin{bmatrix}
    3 \\ 1 \\ 4
\end{bmatrix}\)

Thus,
\[
B = \begin{bmatrix}
\ & \ & \ \\
\vec{b}_1 & \vec{b}_2 & \vec{b}_3 \\
\ & \ & \
\end{bmatrix} = \begin{bmatrix}
    1 & 2 & 3 \\
    1 & 1 & 1 \\
    3 & 4 & 4
\end{bmatrix}
\]

\subsection*{Question 4}
We repeat the process of the last question,

\begin{align*}
    \chi_A(\lambda) = \det \begin{bmatrix}
        -5\lambda & 0 & 4 \\
        -11 & 7 - \lambda & 12 \\
        -1 & 0 & -1 - \lambda
    \end{bmatrix} &= 0
    \\
    = (-5-\lambda)(7-\lambda)(-1-\lambda) + 4(7-\lambda) &= 0
    \\
    \Rightarrow (7-\lambda)(\lambda^2 + 6\lambda + 9) &= 0
\end{align*}
Thus the eigenvalues of \(A\) are \(\lambda = 7\) and \(\lambda = -3\)

To check whether \(A\) is diagonalizable, we need to check the dimension of \(\mathcal{N}(A+3I)\)

\begin{align*}
    \mathcal{N}(A+3I) &, \begin{bmatrix}
        -2 & 0 & 4 &| & 0 \\
        -11 & 10 & 12 &|& 0 \\
        -1 & 0 & 2& |&  0
    \end{bmatrix}
    \\
    &= \begin{bmatrix}
        1 & 0 & -2 &| & 0 \\
        0 & 1 & -1 &|& 0 \\
        0 & 0 & 0& |&  0
    \end{bmatrix}
\end{align*}
Since there is only one linearly independent eigenvector, such as \(\begin{bmatrix}
    2 \\ 1 \\ 1
\end{bmatrix}\), associated with the eigenvalue of \(3\), we conclude that \(A\) is not diagonalizable.
\\

To find a basis of \(\mathbb{R}^3\) consisting of eigenvectors and generalised eigenvectors, we can look at the other eigenvectors 
associated with \(\lambda = 7\) and pick a linearly independent vector from \(\mathcal{N}(A+3I)^2\)

\begin{align*}
    \mathcal{N}(A-7I) &, \begin{bmatrix}
        12 & 0 & 4 &| & 0 \\
        -11 & 0 & 12 &|& 0 \\
        -1 & 0 & -8& |&  0
    \end{bmatrix}
    \\
    &= \begin{bmatrix}
        1 & 0 & 0 &| & 0 \\
        0 & 0 & 1 &|& 0 \\
        0 & 0 & 0& |&  0
    \end{bmatrix}
\end{align*}
One such eigenvector is \(\begin{bmatrix}
    0 \\ 1 \\0
\end{bmatrix}\)

\begin{align*}
    \mathcal{N}(A+3I)^2 &, \begin{bmatrix}
        0 & 0 & 0 &| & 0 \\
        -100 & 100 & 100 &|& 0 \\
        0 & 0 & 0& |&  0
    \end{bmatrix}
    \\
    &= \begin{bmatrix}
        0 & 0 & 0 &| & 0 \\
        -1 & 1 & 1 &|& 0 \\
        0 & 0 & 0& |&  0
    \end{bmatrix}
\end{align*}
A linearly independent vector to the other two eigenvectors in this set is \(\begin{bmatrix}
    3 \\ 1 \\ 2
\end{bmatrix}\)
\\

Hence the set \(\left\{\begin{bmatrix}
0 \\ 1 \\ 0
\end{bmatrix}, \begin{bmatrix}
2 \\ 1 \\1
\end{bmatrix}, \begin{bmatrix}
3 \\ 1 \\ 2
\end{bmatrix}\right\}\) is a basis for \(\mathbb{R}^3\) made up of eigenvectors and a generalised eigenvector

\subsection*{Question 5}
Assume for the sake of contradiction that \(T- \lambda I\) is surjective, 

Then \(\mathrm{Im}(T - \lambda I) = V\)

Also, we know
\begin{align*}
    \dim(\ker(T - \lambda I)) + \dim(\mathrm{Im} \ (T - \lambda I)) &= \dim(V)
    \\
    \Rightarrow \dim(\ker (T - \lambda I)) + \dim(V) &= \dim(V)
    \\
    \Rightarrow \dim(\ker (T - \lambda I)) &= 0
\end{align*}
But since \(\lambda\) is an eigenvalue of \(T\), we know \(T - \lambda I\) has non-zero solutions, since eigenvectors are non-zero, 
this is a contradiction since if the \(\dim(\ker T) = 0\) there exists a unique solution, the trivial one implying there are no non-zero eigenvectors.


\end{document}